% TOP_PROBLEM: {"id": 1204, "title": "The Carathéodory Conjecture", "category_id": 6, "category": {"id": 6, "name": "geometry", "display_name": "Geometry", "description": "Euclidean and non-Euclidean geometry, geometric structures.", "slug": "geometry", "order_index": 6, "created_at": "2026-07-31T15:26:25.671Z"}, "status": "open"}
% FIRST_POSTED: 2026-10-01
% ATTEMPT_STATUS: unresolved
% Imported from: unsolvedmath_additional_500.tex; UnsolvedMath ID 1204
% !TeX program = lualatex
% Compile twice: lualatex 1204.tex
% Self-contained: no external bibliography, images, or \input files.
% Numeric section labels and visible numbers are original UnsolvedMath IDs.
\documentclass[11pt,a4paper]{article}
\usepackage[margin=24mm,headheight=14pt]{geometry}
\usepackage{fontspec}
\setmainfont{Latin Modern Roman}
\setsansfont{Latin Modern Sans}
\setmonofont{Latin Modern Mono}
\usepackage{amsmath,amssymb,mathtools}
\usepackage{unicode-math}
\setmathfont{Latin Modern Math}
\usepackage{microtype}
\usepackage{enumitem,xurl,needspace}
\usepackage[dvipsnames]{xcolor}
\usepackage{hyperref}
\hypersetup{unicode=true,colorlinks=true,linkcolor=MidnightBlue,urlcolor=MidnightBlue,
 pdftitle={UnsolvedMath Problem 1204},pdfauthor={Source-annotated compilation},bookmarksnumbered=true}
\usepackage{bookmark,tocloft,titlesec,fancyhdr}
\setlength{\cftsecnumwidth}{4.2em}
\cftsetpnumwidth{2.5em}
\cftsetrmarg{4em}
\titleformat{\section}{\Large\bfseries\raggedright}{\thesection}{0.7em}{}
\pagestyle{fancy}\fancyhf{}
\fancyhead[L]{\small\sffamily Additional UnsolvedMath statements}
\fancyhead[R]{\small Original catalogue IDs}
\fancyfoot[C]{\thepage}
\setlength{\parindent}{0pt}
\setlength{\parskip}{5pt plus 1pt}
\setlength{\emergencystretch}{3em}
\setcounter{tocdepth}{1}\setcounter{secnumdepth}{1}
\setlist[itemize]{leftmargin=3em,itemsep=4pt,topsep=4pt,parsep=0pt}
\allowdisplaybreaks
\newcommand{\N}{\mathbb{N}}
\newcommand{\Z}{\mathbb{Z}}
\newcommand{\Q}{\mathbb{Q}}
\newcommand{\R}{\mathbb{R}}
\newcommand{\C}{\mathbb{C}}
\newcommand{\F}{\mathbb{F}}
\newcommand{\A}{\mathbb{A}}
\newcommand{\PP}{\mathbb{P}}
\newcommand{\E}{\mathbb{E}}
\newcommand{\eps}{\varepsilon}
\newcommand{\dd}{\,\mathrm d}
\newcommand{\sref}[2]{\hyperlink{source:#1:#2}{[S#2]}}
\newif\ifshowcatalogue
\showcataloguetrue
\newcommand{\cataloguescope}[1]{\ifshowcatalogue
 \paragraph{Statement recorded in the supplied source.}
 #1\fi}
\begin{document}
% UnsolvedMath ID 1204; source catalogue code GEO-013

\setcounter{section}{1203}

\section{Carathéodory's Umbilic-Point Conjecture}\label{1204}

{\small\sffamily UnsolvedMath ID 1204 · Geometry}

\subsection{Definitions and mathematical statement}

\paragraph{Shared notation.}Unless a section says otherwise, $\N=\{1,2,\ldots\}$,
$\Z$ is the ring of integers, $\Q,\R,\C$ are the rational, real, and complex
fields, $[N]=\{1,\ldots,\lfloor N\rfloor\}$, and $\log$ is the natural
logarithm. For real functions of a parameter tending to infinity,
$f\ll g$ and $f=O(g)$ mean $|f|\leq Cg$ eventually for some fixed $C$;
$f\gg g$ reverses this comparison; $f=o(g)$ means $f/g\to0$;
$f\sim g$ means $f/g\to1$; and $f\asymp g$ means both order bounds.
A subscript on $O$ or $\ll$ specifies allowed constant dependence.
The expression $n^{o(1)}$ in an upper bound means growth smaller than
$n^\epsilon$ for each fixed $\epsilon>0$. Local definitions govern any
exception, finite-field convention, or use of the same symbol for another
object. An asymptotic claim is not automatically an effective algorithm.



Let $S\subset\R^3$ be a closed $C^2$ embedded surface bounding a convex body. At a point of $S$, the principal curvatures are the two eigenvalues of its shape operator. A point is umbilic when these eigenvalues agree, equivalently when the second fundamental form is a scalar multiple of the first. The question requires at least two distinct such points. \sref{1204}{1} \sref{1204}{2}

\paragraph{Statement recorded in the supplied source.}
Does every convex, closed, twice-differentiable surface in $\mathbb{R}^3$ have at least two umbilical points? \sref{1204}{1}

\subsection{Short English statement}

Must every smooth convex closed surface have at least two points where its normal curvature is the same in every tangent direction?

\subsection{Research attempt: a topological obstruction and symmetry-forced pairs}
\textbf{Status and regularity audit.} We prove the two-point conclusion for centrally symmetric convex surfaces and for surfaces of revolution. We do not prove it for arbitrary $C^2$ convex surfaces. The inspected abstract of [S2], revised in 2024, asserts a result for $C^{3+\alpha}$ surfaces; that stated regularity does not by itself cover this catalogue's $C^2$ target. No conclusion about the full proof's validity is inferred from its title.

\subsubsection{1. The shape operator must be scalar somewhere}
Because the convex body has interior, choose an interior point as origin. Each ray meets the boundary once; radial projection identifies the boundary topologically with $S^2$. At a $C^2$ point the shape operator $A$ is a continuous self-adjoint endomorphism of the tangent plane. If there were no umbilics, its eigenvalues would be distinct everywhere. Ordering them by size defines a continuous real line subbundle of $TS$.

Here are the precise topological inputs: every real line bundle on $S^2$ is orientable, since its orientation double cover is trivial on the simply connected sphere, and an oriented line bundle admits a continuous nowhere-zero section after a metric normalization. The hairy-ball theorem says $TS^2$ has no such section. Pulling the line bundle back along a smooth radial parametrization gives a contradiction. The parametrization is smooth because the outward normal has strictly positive scalar product with the radial vector from an interior point; the implicit function theorem applies even when a principal curvature is zero. Thus there is at least one umbilic. These standard topological facts, also underlying the index method discussed in [R1], are external inputs; the line-bundle deduction specifies how they enter.

\subsubsection{2. Central symmetry upgrades one point to two}
Suppose $S=-S$ with center the origin. The map $F(x)=-x$ has no fixed point on $S$, since the origin is interior. It is an ambient Euclidean isometry, and the outward normal satisfies $n(-x)=-n(x)$. Differentiating gives the natural conjugacy of shape operators under $dF$; in particular scalarity is preserved. Hence if $x$ is umbilic, so is $-x$, and they are distinct. Together with the previous subsection this proves the conjectured count for every centrally symmetric $C^2$ convex surface, without assuming positive Gaussian curvature.

\subsubsection{3. Rotational symmetry supplies two explicit points}
For a surface invariant under all rotations about an axis through the body's interior, let $p_+,p_-$ be the two intersections of that axis with the boundary. Both are fixed by every rotation. The tangent plane at each is invariant under the rotation group and therefore is the plane perpendicular to the axis. The shape operator commutes with every planar rotation. A self-adjoint real $2\times2$ matrix with that property is a scalar matrix: commuting with rotation through $\pi/2$ already forces its diagonal entries equal and its off-diagonal entry zero. Thus both poles are umbilic. They are distinct because the convex body has nonempty interior.

\subsubsection{4. The index gap in the general case}
For isolated singularities of the principal-direction line field, the line-field Poincare--Hopf formula makes their indices sum to 2 on a sphere. This is a named topological formula, not a proof that each index is at most 1. If such a local upper bound were established under the required geometric hypotheses, one isolated umbilic could not account for the sum. But arbitrary singular line fields can concentrate index 2 at one point; global topology alone excludes zero singularities, not one. Nonisolated umbilics already give at least two points, so the difficult residual case is precisely a single isolated umbilic and a valid local geometric restriction on its index.

The adversarial checks are the fixed-point issue for central inversion, the distinction between convexity and positive curvature, and the regularity of the eigenspace line field. Each has been addressed above. Smoothing a $C^2$ surface and counting two umbilics on approximants is insufficient: two distinct points may converge to the same point. No uniform separation or local index theorem is supplied here. The general target remains unresolved by this attempt.

\subsection{Sources}

{\small

\begin{itemize}

\item[\hypertarget{source:1204:1}{S1}] ULAM.AI, \emph{UnsolvedMath}, supplied \texttt{problems.json}, record 1204 (GEO-013), statement and background. The embedded literature-review date is 2026-08-17; that is the archive’s review date, not a fresh certification. Dataset landing page: \url{https://huggingface.co/datasets/ulamai/UnsolvedMath}.

\item[\hypertarget{source:1204:2}{S2}] Brendan Guilfoyle and Wilhelm Klingenberg, Proof of the Caratheodory Conjecture, arXiv:0808.0851 (2008; revised 2024). \url{https://arxiv.org/abs/0808.0851}. Bibliographic information imported from the supplied record.

\item[\hypertarget{source:1204:3}{S3}] Brendan Guilfoyle and Wilhelm Klingenberg, The three obdurate conjectures of differential geometry, arXiv:2502.11716 (2025). \url{https://arxiv.org/abs/2502.11716}. Bibliographic information imported from the supplied record.

\item[R1] B. Guilfoyle and W. Klingenberg, The three obdurate conjectures of differential geometry (2025); source for the geometric index context.
\url{https://arxiv.org/abs/2502.11716}.
\end{itemize}

}

\label{end:1204}
\end{document}
